\chapter*{Abstract}
\thispagestyle{plain}
\addcontentsline{toc}{chapter}{\numberline{}Abstract}
The brain keeps producing small electrical signals all the time, and these signals change
with what a person is doing and how their mind is working. This project uses a wearable,
four-channel dry-electrode EEG headset to read two kinds of information from the same
recording. The first is which hand a person is imagining moving, left or right, without any
actual movement. This is called motor imagery and it is a basic building block of a
brain-computer interface. The second is the mental state of the person, given as a
cognitive score built from the power in the theta, alpha and beta frequency bands, which
covers attention, cognitive load, relaxation, stress and fatigue.\\

For motor imagery, the analysis method was first checked on a public dataset of 50 acute
stroke patients, and then on the public BCI Competition IV 2b dataset of nine healthy
subjects, which uses the same 250 Hz sampling rate and the same C3 and C4 electrode
positions as the headset. On the 2b data the expected drop in 8--30 Hz power over the side
of the head opposite to the imagined hand was seen in 16 of 18 cases, and the ratio of C3
to C4 power was higher for the left hand than for the right hand in all nine subjects.
Used alone, this ratio identified the hand on unseen days with 73.5\% accuracy, and the
best SVM-based classifier reached 80.2\% on the same person and 74.7\% on people it had
never seen. Left and right motor imagery was also recorded with the headset, where the
same C3/C4 difference appears in the trials examined, but classifiers trained on either
public dataset stayed near chance on the recorded data. This shows that the recording
protocol, rather than the classifier, is the main thing to improve. The cognitive scoring
part works end to end, with the scores computed live against each person's own
calibration. No audio is recorded at any stage.
